\section{Introduction}

Piecewise-constant segmentation---also called multiple change-point inference---is a core
problem in statistics, signal processing, econometrics, computational biology, and machine
learning. In its simplest form, one observes an ordered sequence $y_1,\dots,y_n$ (optionally at
locations $x_1<\cdots<x_n$) and seeks to infer the number of regimes, the boundary locations,
and the segment-specific parameters governing the observations inside each regime. The appeal of
Bayesian formulations is well known: they quantify uncertainty over boundary locations and
segment counts, they support principled model comparison through marginal likelihoods, and they
naturally propagate uncertainty to downstream tasks such as denoising, regression-curve
estimation, and prediction \citep{barry1992ppm, barry1993bayesCP, fearnhead2006exact,
  hutter2006bpcr}.

For offline segmentation, exact Bayesian inference applies whenever single-segment
parameters can be integrated out and the partition prior factorizes appropriately. In that case,
one can precompute block evidences for all candidate segments and then perform dynamic
programming (DP) over contiguous partitions to obtain posterior quantities in polynomial time
\citep{fearnhead2006exact, hutter2006bpcr}. This observation makes piecewise-constant Bayesian
regression unusually modular: once blockwise marginal likelihoods are available, the same global
DP machinery delivers evidences, posterior boundary probabilities, Bayes regression curves, and
joint MAP segmentations.

Despite this favorable structure, the literature remains fragmented. Exact offline Bayesian
methods are often presented for a small number of likelihood/prior pairs, typically in single-
sequence settings, and usually on index-uniform grids. Hierarchical or multi-sequence extensions
are either specialized to particular data types, based on approximate sampling schemes, or focused
on synchrony assumptions that are difficult to combine with a reusable block-evidence interface
\citep{carlin1992hierBayesCP, fearnhead2011dependence, fan2017basic, quinlan2024jrpm}. At the
same time, many modern applications require more than boundary recovery: they need predictive
scores for new sequences, group-level pooling, and transparent uncertainty summaries that can be
exported to downstream decision pipelines.

This paper develops \texttt{BayesBreak}, a unified framework for offline Bayesian segmentation
that separates \emph{local block modeling} from \emph{global partition inference}. The framework is
exact whenever block evidences are available in closed form and extends to controlled deterministic
approximations when the block integral lacks a closed form. The central design principle is simple: if a
candidate segment $(i,j]$ can be summarized by sufficient statistics and assigned a block score
$A^{(0)}_{ij}$ together with any required moment numerators, then the same DP machinery can be
reused across observation models, design priors, and several structured extensions.

\paragraph{Contributions.}
The main contributions are as follows.
\begin{itemize}
\item \textbf{A reusable exponential-family block interface.}
Closed-form block evidences and posterior moments from cumulative sufficient statistics for
exposure-aware EF models with conjugate priors. The construction recovers Gaussian, Poisson, and
Beta--Binomial formulas as special cases.

\item \textbf{Exact offline DP with coherent posterior summaries.}
Forward/backward DP recursions for marginal evidences, segment-count posteriors, boundary
marginals, and Bayes regression curves. These sum-product quantities are distinguished from
the joint MAP segmentation, which uses a separate max-sum recursion with backtracking.

\item \textbf{Design-aware partition priors for irregular designs.}
A segment-length factor (cohesion function) on the partition prior encodes physical spacing
$\Delta_x(t_{q-1},t_q)$ at the prior level, keeping design geometry separate from likelihood
weighting.

\item \textbf{Exact pooling for shared-boundary replicates and known groups.}
Under conditional independence across subjects given shared boundaries, subject-level block evidences multiply, so the pooled posterior is exactly tractable under the same DP; the same construction applies groupwise when memberships are known.

\item \textbf{A latent-group template mixture with exact coordinate updates.}
Each group carries a latent template segmentation; the E-step updates responsibilities and the M-step is an exact responsibility-weighted max-sum DP, yielding a rigorous alternating optimization for a stated finite-template objective rather than full Bayesian averaging.

\item \textbf{Approximate non-conjugate block routines with stability guarantees.}
For non-conjugate GLM blocks we describe Laplace, variational (JJ, P\'olya--Gamma mean field), EP, and 1-D quadrature approximations, and show how a uniform per-block log-evidence error $\varepsilon$ propagates through the DP to perturb global posterior odds by a controlled amount.

\item \textbf{A prediction layer for fitted segmenters.}
Posterior-predictive scoring for new pointwise sequences, set-valued units, and vector-valued responses, distinguishing conditional prediction under an exported segmentation from optional resegmentation-based scoring.
\end{itemize}

\subsection{Related work and literature review}\label{sec:lit}

\paragraph{Frequentist offline segmentation.}
Classical segmentation algorithms optimize additive cost functions over contiguous partitions.
Segment-neighborhood and optimal-partitioning recursions provide exact solutions for a fixed
number of segments or a fixed penalty \citep{auger1989segment, jackson2005optpart}. Later work
reduced practical cost through pruning, most notably in pruned dynamic programming and PELT
\citep{rigaill2010pruned, killick2012pelt}. Convex relaxations such as total variation denoising
and the fused lasso trade exact combinatorial optimization for scalable regularized estimation,
with multisequence extensions such as the group fused lasso encouraging shared boundaries across
signals \citep{rudin1992rof, tibshirani2005fused, bleakley2011groupfused}. Multiscale methods,
including SMUCE and wild binary segmentation, emphasize detection guarantees and confidence
statements across scales rather than full posterior modeling \citep{frick2014smuce,
  fryzlewicz2014wbs}. Circular binary segmentation \citep{olshen2004cbs} is a widely used
recursive detection rule in array-based genomic applications and provides a natural point of
comparison for the Gaussian-block instance of BayesBreak. These methods are indispensable
baselines, but they do not natively return posterior probabilities over partitions or segment
parameters.

\paragraph{Bayesian offline segmentation and product partitions.}
The product-partition formulation of \citet{barry1992ppm, barry1993bayesCP} is one of the most
important starting points for Bayesian change-point analysis. The key idea is that if partition
probabilities factorize over segments and segment parameters can be marginalized, then posterior
inference can be reduced to operations on block scores. \citet{fearnhead2006exact} gave an exact
and efficient DP-based treatment for a broad class of offline changepoint models, while
\citet{hutter2006bpcr} showed how Bayesian averaging over partitions yields regression curves,
boundary probabilities, and evidence-based model selection in the piecewise-constant setting.
Earlier work by \citet{yao1988biometrika} studied changepoint estimation via Schwarz's criterion,
complementing the evidence-based Bayesian model-selection view. The astronomical \emph{Bayesian
Blocks} construction of \citet{scargle2013bayesianblocks} is essentially a Poisson-block instance
of the product-partition program and illustrates how domain-specific block routines combine with a
single DP recursion. Work by \citet{fearnhead2011dependence} and \citet{ruggieri2014bcpvs}
demonstrated how far the DP viewpoint can be pushed when one adds dependence across segments or
richer within-segment regressions. \citet{punskaya2002bayesian} applied a related MCMC-based
Bayesian curve-fitting program to noisy signals.

\paragraph{Reviews and surveys.}
A selective survey of modern offline change-point methods, including penalized, Bayesian,
nonparametric, and multiscale variants, is given by \citet{truong2020review}; we view BayesBreak
as sitting inside the Bayesian exact-offline cell of that taxonomy, broadened across block models.

\paragraph{Bayesian online and approximate methods.}
When streaming updates are required, Bayesian online changepoint detection provides an online
filtering alternative that computes an exact run-length posterior rather than an exact offline
partition posterior \citep{adamsmackay2007bocpd}. When exact
marginalization is unavailable, reversible-jump MCMC and related sampling methods remain central
Bayesian tools \citep{green1995rjMCMC, denison1998bayesian}. BayesBreak is aimed squarely at the
offline regime: we prioritize reusable block marginalization and exact DP whenever the model
permits it, and use deterministic local approximations only when the block integral itself lacks a
closed form.

\paragraph{Multi-sequence and structured Bayesian segmentation.}
Many applications exhibit boundary sharing or boundary dependence across related sequences.
Hierarchical Bayesian changepoint models appear in early work by \citet{carlin1992hierBayesCP},
while more recent multisequence methods such as BASIC exploit cross-sequence co-occurrence of
changepoints through hierarchical priors \citep{fan2017basic}. Joint random partition models
provide a broader Bayesian nonparametric framework for borrowing information across multiple
ordered processes \citep{quinlan2024jrpm}. Related product-partition literature also studies
covariate-aware random partitions and generalized cohesion functions, which are conceptually close
to the design-aware partition priors used here \citep{muller2011ppmx, park2010gppm}.

\paragraph{Positioning of BayesBreak.}
BayesBreak occupies a modular middle ground in this landscape: it retains exact offline DP for the core single-sequence and
shared-boundary pooled settings, broadens the menu of admissible block models, clarifies how to
incorporate irregular designs through the prior rather than through ad hoc weighting, and exports a
prediction interface that is usually missing from segmentation papers. For latent groups, BayesBreak
uses a template-mixture extension whose EM updates are exact for the stated finite objective.

\paragraph{Annotated review.}
A compact annotated literature review, organized by methodological theme and relation to
BayesBreak, is provided in Appendix~\ref{app:annotated-lit}.

\paragraph{Code and reproduction.}
A reference implementation of BayesBreak, including the block routines, DP layer, and the
reproduction pipelines for the four real-data case studies, is available at
\url{https://github.com/osolari/bayesbreak}. Installation, version pinning, and end-to-end
commands are collected in Appendix~\ref{app:code}.

\paragraph{Paper organization.}
The paper is in three parts. Theory (\S\ref{sec:problem}--\ref{sec:families}): inferential
targets, notation, the Bayesian model, the EF--conjugate block layer, the DP, irregular designs,
replicate and group pooling, the latent-group template mixture, and family-specific block
formulas. Non-conjugate methods and prediction (\S\ref{sec:nonconj}--\ref{sec:prediction}):
deterministic block approximations with stability theory, and the posterior-predictive interface
for new pointwise, set-valued, and vector-valued data. Algorithms, limitations, experiments, and
conclusions (\S\ref{sec:algorithms}--\ref{sec:conclusion}).
